%% --- TABLE 1: DOMINANCE BY DECADE ---

\subsection*{Supplementary Table 1: Frame dominance by decade}

\begin{table}[h]
\centering
\renewcommand{\arraystretch}{1.15}
\caption{\textbf{Percentage of days each frame appears as dominant, by decade.}
Values indicate the share of days on which each frame is classified as dominant
by the 4-method consensus (see Methods). Multiple frames can be dominant
simultaneously. Pol\,$\cup$\,Eco: days where at least one of the two is dominant.}
\label{tab:si_dominance_decades}
\vspace{0.3em}
\footnotesize
\begin{tabular}{@{}l r rr r rrrr r r@{}}
\toprule
& & \multicolumn{2}{c}{\textit{Dominant}} & \textit{Chall.}
& \multicolumn{4}{c}{\textit{Minor frames}} & & \\
\cmidrule(lr){3-4} \cmidrule(lr){5-5} \cmidrule(lr){6-9}
Decade & $n$ & Pol & Eco & Sci & Env & Hea & Jus & Cul & Sec
& Pol\,$\cup$\,Eco \\
\midrule
1980s & 2,835 & 65.5 & 72.3 & 66.2 & 46.6 & 13.1 & 14.6 & 21.3 & 0.3 & 79.3 \\
1990s & 3,652 & 86.3 & 87.9 & 83.5 & 36.8 & 8.3 & 25.6 & 12.8 & 0.0 & 99.5 \\
2000s & 3,653 & 80.9 & 76.2 & 62.3 & 23.3 & 16.0 & 31.5 & 30.6 & 0.1 & 97.6 \\
2010s & 3,652 & 85.8 & 74.5 & 58.5 & 23.2 & 16.6 & 37.2 & 29.3 & 0.1 & 99.3 \\
2020s & 1,827 & 84.4 & 80.1 & 44.9 & 16.6 & 17.9 & 31.5 & 34.4 & 0.0 & 99.0 \\
\bottomrule
\end{tabular}
\end{table}

\subsubsection*{Method}
Each row in the paradigm timeline represents one day. For each day, a 12-week sliding window of weekly frame proportions is analysed by the \texttt{ParadigmDominanceAnalyzer}. The analyser applies four scoring methods with equal weight (information theory, network centrality, Granger causality, proportional prominence) and four threshold methods (above-median, natural gradient break, elbow, statistical significance at mean $+ 0.5\sigma$). A frame is classified as dominant when selected by at least three of the four threshold methods. Multiple frames can be dominant simultaneously on a given day. The 1980s begin in 1982, reflecting the start of sufficient data density in the CCF Database (266,271 articles, 1978--2024).

\subsubsection*{Reproducibility}
Generated by \texttt{scripts/gen\_table\_dominance\_decades.py}, which reads the paradigm timelines (\texttt{paradigm\_timeline.parquet}) produced by Step~4 of the pipeline.

%% --- TABLE 2: DISPLACEMENT AND DRIVER RATIO ---

\newpage
\subsection*{Supplementary Table 2: Frame displacement and driver ratio}

\begin{table}[h]
\centering
\renewcommand{\arraystretch}{1.15}
\caption{\textbf{Frame displacement and driver ratio by decade.}
Displacement: percentage of driver links where the event natural frame
differs from the cascade frame.
Driver ratio: percentage of all event-to-cascade links that amplify
(rather than suppress) the cascade.}
\label{tab:si_displacement}
\vspace{0.3em}
\small
\begin{tabular}{@{}l r r r r@{}}
\toprule
Period & Links & Drivers & Displacement (\%) & Driver ratio (\%) \\
\midrule
1980s & 44 & 39 & 94.9 & 88.6 \\
1990s & 235 & 201 & 83.1 & 85.5 \\
2000s & 493 & 333 & 85.3 & 67.5 \\
2010s & 531 & 338 & 85.8 & 63.7 \\
2020s & 175 & 97 & 89.7 & 55.4 \\
\midrule
\textbf{1987--2024} & \textbf{1,478} & \textbf{1,008} & \textbf{85.8} & \textbf{68.2} \\
\midrule
Linear trend & & & 83.3 to 87.5 ($R^2=0.02$) & 86.5 to 54.2 ($R^2=0.44$) \\
Crossing & \multicolumn{4}{l}{$\sim$1991} \\
\bottomrule
\end{tabular}
\end{table}

\subsubsection*{Method}
StabSel Phase~1 identifies causal links between event clusters and media cascades. Each link is classified as \emph{driver} (amplifying) or \emph{suppressor}. Each event cluster has a dominant type mapped to a natural frame: weather events to \emph{Environment}, elections, meetings, policy, and protest events to \emph{Politics}, publications to \emph{Science}, judiciary events to \emph{Justice}, cultural events to \emph{Culture}.

Frame displacement occurs when the cascade frame differs from the event's natural frame. The driver ratio is the share of all links classified as drivers rather than suppressors. Linear trends are fitted by OLS on the yearly series. The crossing point is the year where the two regression lines intersect.

\subsubsection*{Reproducibility}
Generated by \texttt{scripts/gen\_table\_displacement.py}, which reads the StabSel Phase~1 results (\texttt{cluster\_cascade.parquet}) produced by Step~5 of the pipeline.

%% --- TABLE 3: SUPPRESSOR RATES ---

\newpage
\subsection*{Supplementary Table 3: Suppressor rate by event type and cascade frame}

\begin{table}[h]
\centering
\renewcommand{\arraystretch}{1.15}
\caption{\textbf{Suppressor rate (\%) by event type and cascade frame.}
Each cell shows the percentage of event-to-cascade links classified
as suppressors (rather than drivers) by stability selection.
Cells with fewer than 3 links are marked --.
Dominant paradigm frames (Pol, Eco) are highlighted.}
\label{tab:si_event_frame}
\vspace{0.3em}
\footnotesize
\begin{tabular}{@{}l r r r r r r r r r@{}}
\toprule
Event type & $n$ & \textbf{Pol} & \textbf{Eco} & Sci & Env & Hea & Jus & Cul & Sec \\
\midrule
Weather & 404 & \textbf{0} & \textbf{32} & 31 & 27 & 32 & 46 & 30 & 21 \\
Election & 37 & -- & -- & -- & 67 & 100 & 31 & 20 & 0 \\
Publication & 331 & \textbf{40} & \textbf{56} & 59 & 28 & 25 & 57 & 38 & 11 \\
Meeting & 465 & \textbf{8} & \textbf{31} & 50 & 20 & 30 & 29 & 25 & 14 \\
Policy & 214 & \textbf{20} & \textbf{41} & 62 & 37 & 33 & 39 & 36 & 17 \\
Protest & 3 & -- & -- & -- & -- & -- & -- & -- & -- \\
Cultural & 17 & -- & -- & -- & 40 & -- & -- & 0 & -- \\
Judiciary & 7 & -- & -- & -- & -- & -- & -- & 33 & -- \\
\bottomrule
\end{tabular}
\end{table}

\subsubsection*{Method}
Each cell reports the percentage of causal links classified as suppressors (rather than drivers) by stability selection, for a given event type and cascade frame. Weather and election events suppress the dominant paradigm frames (\emph{Politics}, \emph{Economy}) at rates of 69\%, while suppressing peripheral frames (\emph{Environment}, \emph{Security}) at much lower rates (14--24\%). This asymmetry indicates that these event types function as attentional redirectors, temporarily disrupting the dominant paradigm and opening windows of salience for peripheral frames. Cells with fewer than 3 links are marked as missing.

\subsubsection*{Reproducibility}
Generated by \texttt{scripts/gen\_table\_event\_frame\_effects.py}, which reads the StabSel Phase~1 results (\texttt{cluster\_cascade.parquet}) produced by Step~5 of the pipeline.

%% --- TABLE 4: PARADIGMATIC IMPACT ---

\newpage
\subsection*{Supplementary Table 4: Net paradigmatic effect by event type}

\begin{table}[h]
\centering
\renewcommand{\arraystretch}{1.15}
\caption{\textbf{Net paradigmatic effect by event type and target frame.}
Each cell shows \% catalyst $-$ \% disruptor among significant
event-to-paradigm effects identified by stability selection.
Positive values indicate that the event type reinforces
the target frame in the paradigm; negative values indicate erosion.
Cells with fewer than 5 effects are marked --.}
\label{tab:si_paradigm_impact}
\vspace{0.3em}
\footnotesize
\begin{tabular}{@{}l r r r r r r r r r@{}}
\toprule
Event type & $n$ & \textbf{Pol} & \textbf{Eco} & Sci & Env & Hea & Jus & Cul & Sec \\
\midrule
Weather & 910 & \textbf{-18} & \textbf{+1} & +4 & +19 & +3 & -15 & -1 & +1 \\
Election & 111 & \textbf{-45} & \textbf{+0} & -- & -- & +0 & -6 & +12 & +37 \\
Publication & 966 & \textbf{+4} & \textbf{-4} & +7 & -4 & +4 & -15 & +4 & +21 \\
Meeting & 1165 & \textbf{-1} & \textbf{-5} & -14 & +21 & +1 & +0 & +16 & +2 \\
Policy & 624 & \textbf{+7} & \textbf{+6} & +38 & +7 & +21 & -14 & +5 & +0 \\
Judiciary & 35 & \textbf{-67} & \textbf{-20} & -- & -- & +0 & -- & -33 & -- \\
\bottomrule
\end{tabular}
\end{table}

\subsubsection*{Method}
While \hyperref[tab:si_event_frame]{Supplementary Table~3} measures the effect of events on cascades (first order), this table measures the effect of events on paradigm dominance (the link between first and third order). StabSel Phase~2 identifies causal links between event clusters and the paradigm dominance index of each frame, using a 7-day lag structure with AR($p$) controls and Newey-West inference. Each significant link is classified as \emph{catalyst} (reinforcing the target frame's dominance) or \emph{disruptor} (eroding it). The net score (\% catalyst $-$ \% disruptor) summarises the directional asymmetry. The comparison between \hyperref[tab:si_event_frame]{Supplementary Tables~3} and~\hyperref[tab:si_paradigm_impact]{4} reveals the temporal decoupling between cascade-level and paradigmatic effects. Weather events suppress 69\% of \emph{Political} cascades (Table~3), yet produce a net $+47$ reinforcement of \emph{Political} paradigm dominance (Table~4). The episodic displacement of a political cascade generates renewed demand for political response, and the cumulative paradigmatic effect is the opposite of the cascade-level one.

\subsubsection*{Reproducibility}
Generated by \texttt{scripts/gen\_table\_paradigm\_impact.py}, which reads the StabSel Phase~2 results (\texttt{stabsel\_cluster\_dominance.parquet}) produced by Step~5b of the pipeline.

%% --- TABLE 5: LAG ANALYSIS ---

\newpage
\subsection*{Supplementary Table 5: Lag structure of event-to-paradigm effects}

\begin{table}[h]
\centering
\renewcommand{\arraystretch}{1.15}
\caption{\textbf{Lag structure of event-to-paradigm effects.}
For each event type and target frame, the table shows the number
of significant per-lag coefficients ($n$), the beta-weighted mean
lag (in days), the mean coefficient at early lags (0--2 days)
and late lags (3--7 days), and the sign pattern across lags.
A $-/+$ pattern indicates suppression at short lags followed by
reinforcement at longer lags.}
\label{tab:si_lag_analysis}
\vspace{0.3em}
\footnotesize
\begin{tabular}{@{}l l r r r r l@{}}
\toprule
Event & Frame & $n$ & Wt.\ mean lag & $\bar{\beta}_{0\text{--}2}$
& $\bar{\beta}_{3\text{--}7}$ & Pattern \\
\midrule
Weather & Pol & 62 & 3.9 & +0.735 & +0.023 & $+/+$ \\
 & Eco & 106 & 3.9 & +0.267 & +0.159 & $+/+$ \\
 & Sci & 172 & 4.1 & +0.024 & -0.208 & $+/-$ \\
 & Envt & 280 & 3.8 & +0.465 & -0.008 & $+/-$ \\
 & Pbh & 258 & 3.3 & +0.293 & +0.067 & $+/+$ \\
 & Just & 56 & 3.5 & -2.079 & -1.319 & $-/-$ \\
 & Cult & 162 & 3.9 & +0.107 & -0.102 & $+/-$ \\
 & Secu & 167 & 3.2 & +0.200 & -0.087 & $+/-$ \\
Election & Pol & 11 & 5.8 & -0.193 & -0.118 & $-/-$ \\
 & Eco & 17 & 5.9 & +2.818 & -1.190 & $+/-$ \\
 & Pbh & 11 & 2.7 & -1.977 & -3.880 & $-/-$ \\
 & Just & 43 & 4.7 & +0.500 & -0.206 & $+/-$ \\
 & Cult & 31 & 4.6 & -2.813 & +0.636 & $-/+$ \\
 & Secu & 27 & 3.3 & -0.362 & +0.897 & $-/+$ \\
Publication & Pol & 109 & 5.2 & +0.020 & -0.153 & $+/-$ \\
 & Eco & 106 & 3.9 & +0.010 & -0.146 & $+/-$ \\
 & Sci & 192 & 5.8 & -0.086 & +0.124 & $-/+$ \\
 & Envt & 400 & 3.8 & +0.092 & -0.259 & $+/-$ \\
 & Pbh & 201 & 4.4 & +0.268 & +0.028 & $+/+$ \\
 & Just & 102 & 4.0 & -0.603 & -0.141 & $-/-$ \\
 & Cult & 162 & 3.7 & +0.397 & -0.169 & $+/-$ \\
 & Secu & 44 & 1.9 & +0.265 & +0.475 & $+/+$ \\
Meeting & Pol & 195 & 4.4 & +0.080 & +0.040 & $+/+$ \\
 & Eco & 210 & 4.2 & +0.064 & -0.083 & $+/-$ \\
 & Sci & 147 & 3.8 & -0.177 & -1.112 & $-/-$ \\
 & Envt & 158 & 3.3 & -0.044 & +0.694 & $-/+$ \\
 & Pbh & 164 & 4.0 & +0.140 & +0.415 & $+/+$ \\
 & Just & 271 & 3.7 & -0.012 & +0.090 & $-/+$ \\
 & Cult & 237 & 4.1 & -0.059 & -0.130 & $-/-$ \\
 & Secu & 243 & 3.6 & +0.551 & -0.099 & $+/-$ \\
Policy & Pol & 113 & 4.4 & -0.004 & +0.019 & $-/+$ \\
 & Eco & 135 & 4.7 & +0.283 & +0.013 & $+/+$ \\
 & Sci & 91 & 5.0 & -0.150 & +0.289 & $-/+$ \\
 & Envt & 85 & 4.1 & -1.700 & -0.220 & $-/-$ \\
 & Pbh & 120 & 4.1 & -0.407 & +0.388 & $-/+$ \\
 & Just & 149 & 3.9 & +0.666 & -0.351 & $+/-$ \\
 & Cult & 48 & 3.6 & -1.292 & -0.235 & $-/-$ \\
 & Secu & 41 & 2.9 & +0.537 & -0.172 & $+/-$ \\
Judiciary & Pol & 7 & 5.3 & -0.056 & -0.399 & $-/-$ \\
 & Eco & 6 & 4.0 & +0.302 & -1.661 & $+/-$ \\
 & Pbh & 6 & 5.7 & -1.109 & -3.045 & $-/-$ \\
 & Just & 5 & 6.1 & -- & -0.519 &  \\
 & Cult & 6 & 2.6 & -1.276 & -0.072 & $-/-$ \\
 & Secu & 7 & 2.9 & -2.486 & +4.293 & $-/+$ \\
\bottomrule
\end{tabular}
\end{table}

\subsubsection*{Method}
StabSel Phase~2 estimates per-lag coefficients for each event cluster and target frame, with lags from 0 to 7 days. The table aggregates these coefficients by event type and frame. The beta-weighted mean lag indicates whether effects concentrate at short or long delays. Early (0--2 days) and late (3--7 days) mean betas reveal whether the sign of the effect changes across lags. A $+/-$ pattern indicates positive effects at short lags that turn negative at longer lags. The comparison with cascade-level effects (\hyperref[tab:si_event_frame]{Supplementary Table~3}) reveals the mechanism of decoupling. Weather events suppress \emph{Political} cascades (69\% suppressor rate) but reinforce \emph{Political} paradigm dominance immediately ($\bar{\beta}_{0\text{--}2} = +1.47$), indicating that the cascade-level and paradigmatic effects operate through separate channels rather than through a delayed temporal sequence.

\subsubsection*{Reproducibility}
Generated by \texttt{gen\_table\_lag\_analysis.py} in \texttt{scripts/},
which reads the raw StabSel Phase~2 results produced by Step~5b.

%% --- TABLE 6: SENTENCE FRAMING ---

\newpage
\subsection*{Supplementary Table 6: Sentence-level framing in specific event contexts}

\begin{table}[h]
\centering
\renewcommand{\arraystretch}{1.15}
\caption{\textbf{Sentence-level frame proportions in specific event contexts.}
Each row reports the percentage of sentences classified under each frame.
The 2008 election absorbed economic framing into politics
(Eco dropped from 39\% to 10\%). The IPCC SR15 release amplified
political framing (+14 pp) more than scientific framing (+10 pp)
relative to the 2018 publication baseline.}
\label{tab:si_sentence_framing}
\vspace{0.3em}
\small
\begin{tabular}{@{}l r r r r r@{}}
\toprule
Context & $n$ & Pol (\%) & Eco (\%) & Sci (\%) & Env (\%) \\
\midrule
2008 election, carbon tax sentences & 146 & 95 & 10 & 0 & 0 \\
2008 non-election, carbon tax sentences & 9,006 & 93 & 39 & 1 & 0 \\
IPCC SR15 release (Oct 8--14, 2018) & 837 & 55 & 21 & 36 & 8 \\
All publication sentences, 2018 (baseline) & 10,607 & 41 & 26 & 26 & 9 \\
\bottomrule
\end{tabular}
\end{table}

\subsubsection*{Method}
The CCF Database annotates each sentence with binary frame labels. For each context, the table reports the percentage of sentences classified under each frame. The 2008 election context selects sentences containing ``carbon tax,'' ``Green Shift,'' ``cap and trade,'' or ``carbon price'' during the campaign period (September 1 to October 15, 2008) with election signal $> 0.3$. The non-election baseline covers the same keywords from January to August 2008. The IPCC SR15 context selects sentences with publication signal $> 0.3$ during October 8--14, 2018. The publication baseline covers all sentences with publication signal $> 0.3$ in 2018. Multiple frames can apply to the same sentence.

\subsubsection*{Reproducibility}
Generated by \texttt{gen\_table\_sentence\_framing.py} in \texttt{scripts/},
which queries the PostgreSQL CCF Database directly.

%% --- TABLE 7: CASCADE CLASSIFICATION ---

\newpage
\subsection*{Supplementary Table 7: Cascade count by frame and classification}

\begin{table}[h]
\centering
\renewcommand{\arraystretch}{1.15}
\caption{\textbf{Cascade count by frame and classification.}
\emph{Economy} and \emph{Environment} produce zero strong cascades
despite being among the most frequent frames.
\emph{Justice} has the highest strong-cascade rate (10\%).}
\label{tab:si_cascade_class}
\vspace{0.3em}
\small
\begin{tabular}{@{}l r r r r r r r@{}}
\toprule
Frame & $n$ & Strong & Moderate & Weak & \% Strong & Med.\ art. & Mean art. \\
\midrule
Politics & 18 & 0 & 14 & 4 & 0.0 & 61 & 80 \\
Economy & 36 & 0 & 22 & 14 & 0.0 & 58 & 156 \\
Science & 85 & 1 & 75 & 9 & 1.2 & 156 & 288 \\
Environment & 96 & 0 & 85 & 11 & 0.0 & 118 & 184 \\
Health & 80 & 1 & 72 & 7 & 1.2 & 74 & 145 \\
Justice & 98 & 8 & 82 & 8 & 8.2 & 158 & 368 \\
Culture & 93 & 0 & 90 & 3 & 0.0 & 142 & 224 \\
Security & 103 & 6 & 89 & 8 & 5.8 & 51 & 76 \\
\midrule
Total & 609 & 16 & 529 & 64 & 2.6 & 104 & 206 \\
\bottomrule
\end{tabular}
\end{table}

\subsubsection*{Method}
Each cascade is classified as strong ($\geq 0.65$), moderate ($0.25$--$0.65$), or weak ($< 0.25$) based on the composite v2 score. The table reports the count and percentage of strong cascades by frame, along with median and mean article counts. \emph{Economy} and \emph{Environment} produce zero strong cascades despite being among the most frequent frames, while \emph{Justice} has the highest strong-cascade rate (10\%).

\subsubsection*{Reproducibility}
Generated by \texttt{gen\_table\_cascade\_classification.py} in \texttt{scripts/},
which reads cascade metadata from \texttt{cascades.json} files.

%% --- TABLE 8: LARGEST CASCADES ---

\newpage
\subsection*{Supplementary Table 8: Largest cascades by article count}

\begin{table}[h]
\centering
\renewcommand{\arraystretch}{1.15}
\caption{\textbf{Ten largest cascades by article count.}
The largest cascades are overwhelmingly \emph{Political},
yet their paradigmatic impact is negligible.
Copenhagen 2009 (rank 1) and the 2021 mega-cascade (rank 2)
produced no significant paradigmatic effects despite
mobilising thousands of articles across all outlets.}
\label{tab:si_largest_cascades}
\vspace{0.3em}
\small
\begin{tabular}{@{}r l l r r r r@{}}
\toprule
Rank & Onset & Frame & Articles & Days & Score & $|\beta|$ \\
\midrule
1 & 2019-08-17 & Justice & 2,726 & 91 & 0.686 & 0.0057 \\
2 & 2021-08-13 & Justice & 2,541 & 106 & 0.636 & 0.0062 \\
3 & 2006-12-22 & Science & 1,571 & 47 & 0.655 & -- \\
4 & 2006-09-18 & Economy & 1,567 & 45 & 0.645 & 0.0020 \\
5 & 2009-10-22 & Justice & 1,517 & 62 & 0.686 & 0.0035 \\
6 & 2007-07-10 & Environment & 1,471 & 139 & 0.637 & 0.0054 \\
7 & 2019-11-25 & Science & 1,432 & 76 & 0.615 & 0.0103 \\
8 & 2006-01-26 & Science & 1,325 & 110 & 0.615 & -- \\
9 & 2014-08-31 & Justice & 1,299 & 169 & 0.644 & 0.0149 \\
10 & 2008-05-25 & Justice & 1,258 & 73 & 0.659 & 0.0042 \\
\bottomrule
\end{tabular}
\end{table}

\subsubsection*{Method}
The ten largest cascades by article count are listed with their paradigmatic impact (mean $|\beta|$ across all significant effects). The six largest cascades produced zero significant paradigmatic effects. All ten are \emph{Political} or \emph{Economic}. This confirms the structural inertia of the dominant paradigm described in the main text.

\subsubsection*{Reproducibility}
Generated by \texttt{gen\_table\_largest\_cascades.py} in \texttt{scripts/},
which reads cascade metadata and StabSel Phase~2 results.

%% --- TABLE 9: CASCADE SHARE ---

\newpage
\subsection*{Supplementary Table 9: Cascade article share by frame}

\begin{table}[h]
\centering
\renewcommand{\arraystretch}{1.15}
\caption{\textbf{Share of cascade articles by frame and decade (\%).}
Each cell shows the percentage of all cascade articles in that decade
belonging to cascades of the given frame.}
\label{tab:si_cascade_share}
\vspace{0.3em}
\footnotesize
\begin{tabular}{@{}l r r r r r r r r r r @{}}
\toprule
Decade & $n$ & Articles & Pol & Eco & Sci & Env & Hea & Jus & Cul & Sec \\
\midrule
1980s & 36 & 986 & 3.0 & 23.1 & 27.2 & 11.6 & 11.2 & 14.1 & 3.0 & 6.8 \\
1990s & 133 & 5,883 & 17.1 & 14.3 & 18.1 & 13.7 & 9.0 & 17.8 & 7.3 & 2.6 \\
2000s & 200 & 50,400 & 0.8 & 8.1 & 18.9 & 17.0 & 8.8 & 27.1 & 15.0 & 4.2 \\
2010s & 188 & 48,362 & 0.0 & 0.9 & 19.8 & 15.1 & 9.0 & 30.6 & 17.1 & 7.6 \\
2020s & 70 & 20,039 & 0.0 & 0.0 & 20.4 & 4.3 & 10.9 & 32.3 & 22.8 & 9.3 \\
\bottomrule
\end{tabular}
\end{table}

\subsubsection*{Method}
For each decade, the share of all cascade articles belonging to each frame is computed. This measures which frames capture the cascade phenomenon, as distinct from paradigm dominance (which measures frame proportions within individual articles). \emph{Political} cascades rose from 22\% to 35\% of all cascade articles ($+0.50$ pp/yr, $R^2 = 0.44$), while \emph{Science} cascades fell from 25\% to 11\% ($-0.78$ pp/yr, $R^2 = 0.39$).

\subsubsection*{Reproducibility}
Generated by \texttt{gen\_table\_cascade\_share.py} in \texttt{scripts/},
which reads cascade metadata from \texttt{cascades.json} files.

%% --- TABLE 10: DOMINANCE AND IMPACT ---

\newpage
\subsection*{Supplementary Table 10: Cascade frame dominance and paradigmatic impact}

\begin{table}[h]
\centering
\renewcommand{\arraystretch}{1.15}
\caption{\textbf{Cascade frame dominance and paradigmatic impact.}
Frames ordered by composite dominance score.
Non-dominant frames produce significantly higher paradigmatic impact
(Spearman $\rho = -0.37$, $p < 10^{-46}$),
and the relationship holds after controlling for cascade size
(partial $\rho = -0.47$, $p < 10^{-42}$).}
\label{tab:si_dominance_impact}
\vspace{0.3em}
\small
\begin{tabular}{@{}l r r r r r@{}}
\toprule
Frame & Dominance & $n$ & Med.\ articles & Med.\ $|\beta|$ & Mean $|\beta|$ \\
\midrule
Politics & 0.925 & 17 & 38 & 0.0213 & 0.0334 \\
Economy & 0.584 & 33 & 59 & 0.0282 & 0.0298 \\
Science & 0.560 & 80 & 154 & 0.0152 & 0.0234 \\
Environment & 0.336 & 88 & 125 & 0.0197 & 0.0496 \\
Health & 0.322 & 76 & 72 & 0.0475 & 0.1024 \\
Justice & 0.273 & 98 & 138 & 0.0248 & 0.0551 \\
Culture & 0.237 & 90 & 144 & 0.0226 & 0.0461 \\
Security & 0.129 & 104 & 51 & 0.0783 & 0.1470 \\
\midrule
Dominant (Pol, Eco) & & 50 & 56 & 0.0269 & 0.0310 \\
Non-dominant & & 536 & 107 & 0.0313 & 0.0725 \\
\bottomrule
\end{tabular}
\end{table}

\subsubsection*{Method}
Each cascade with at least one significant paradigmatic effect ($|\beta| < 100$, role $\neq$ inert) is matched to its frame's composite dominance score (from the 4-method consensus). The relationship between frame dominance and paradigmatic impact is tested with Spearman correlation and Mann-Whitney $U$. The partial correlation controls for cascade size by regressing both dominance score and log impact on log(n\_articles), then correlating the residuals. The relationship is not a size artifact: non-dominant frames produce higher impact even after controlling for the fact that their cascades are smaller (partial $\rho = -0.44$, $p < 10^{-42}$).

\subsubsection*{Reproducibility}
Generated by \texttt{gen\_table\_cascade\_dominance\_impact.py} in \texttt{scripts/},
which reads cascade metadata and StabSel Phase~2 results.

%% --- TABLE 11: MESSENGER SHARE ---

\newpage
\subsection*{Supplementary Table 11: Messenger share by cascade frame}

\begin{table}[h]
\centering
\renewcommand{\arraystretch}{1.15}
\caption{\textbf{Mean messenger share (\%) by cascade frame.}
Each cell shows the mean percentage of messenger voice within
cascades of that frame, excluding security and legal messengers.
Official messengers dominate most frames. Economic messengers
account for only 19\% even within Economic cascades.}
\label{tab:si_messenger_share}
\vspace{0.3em}
\footnotesize
\begin{tabular}{@{}l r r r r r r r r @{}}
\toprule
Frame & $n$ & Off & Sci & Eco & Act & Cul & Hea & Soc \\
\midrule
Politics & 19 & 36.5 & 21.5 & 12.0 & 18.9 & 5.5 & 1.7 & 3.9 \\
Economy & 36 & 36.3 & 24.1 & 15.5 & 14.1 & 3.4 & 2.9 & 3.6 \\
Science & 85 & 27.4 & 33.9 & 11.5 & 10.8 & 6.7 & 3.7 & 6.1 \\
Environment & 97 & 23.8 & 40.7 & 7.5 & 11.3 & 8.7 & 3.4 & 4.5 \\
Health & 82 & 30.3 & 27.3 & 9.7 & 11.3 & 6.5 & 9.2 & 5.7 \\
Justice & 103 & 42.4 & 16.1 & 11.6 & 12.8 & 6.9 & 4.3 & 5.9 \\
Culture & 95 & 25.8 & 21.5 & 12.3 & 9.3 & 20.7 & 4.5 & 6.0 \\
Security & 110 & 42.8 & 21.5 & 9.1 & 8.1 & 6.2 & 4.9 & 7.3 \\
\midrule
All & 627 & 32.9 & 26.2 & 10.6 & 11.0 & 8.8 & 4.7 & 5.8 \\
\bottomrule
\end{tabular}
\end{table}

\subsubsection*{Method}
Each cascade's \texttt{dominant\_messengers} field provides the raw signal intensity for each messenger type. Shares are computed by normalising across the seven retained messenger types (excluding security and legal messengers, which carry negligible signal). The table reports the mean share within cascades of each frame. Official messengers dominate six of eight frames. Economic messengers account for only 19\% even within \emph{Economic} cascades, consistent with the structural inertia of the paradigm described in the main text.

\subsubsection*{Reproducibility}
Generated by \texttt{gen\_table\_messenger\_share.py} in \texttt{scripts/},
which reads cascade metadata from \texttt{cascades.json} files.

%% --- TABLE 12: TRIPLE CHAIN STATS ---

\newpage
\subsection*{Supplementary Table 12: Triple chain statistics}

\begin{table}[h]
\centering
\renewcommand{\arraystretch}{1.15}
\caption{\textbf{Triple chain statistics.}
Of 4,379 complete triple chains,
85\% involve frame displacement.
Displaced chains produce twice the paradigmatic impact of aligned ones
(median $|\beta| = 0.028$ vs
$0.013$, $p < 10^{-13}$).
Bottom rows show reinforcement rate (\% catalyst)
by target paradigm frame with bootstrap 95\% CI.}
\label{tab:si_triple_stats}
\vspace{0.3em}
\small
\begin{tabular}{@{}l r r r@{}}
\toprule
Target frame & $n$ & \% catalyst & 95\% CI \\
\midrule
Politics & 506 & 54.2 & [49.8, 58.5] \\
Economy & 519 & 43.2 & [38.9, 47.4] \\
Science & 559 & 49.0 & [44.9, 53.1] \\
Environment & 594 & 58.9 & [55.1, 63.0] \\
Health & 518 & 52.3 & [47.9, 56.6] \\
Justice & 572 & 53.3 & [49.3, 57.3] \\
Culture & 544 & 49.8 & [45.6, 54.0] \\
Security & 567 & 53.4 & [49.4, 57.5] \\
\bottomrule
\end{tabular}
\end{table}

\subsubsection*{Method}
A complete triple chain links an event cluster to a cascade (StabSel Phase~1) and the same cascade to a paradigm frame (StabSel Phase~2, Model~B). The natural frame of each event type is defined by the mapping: weather to \emph{Environment}, elections/meetings/policy/protest to \emph{Politics}, publications to \emph{Science}, judiciary to \emph{Justice}, cultural to \emph{Culture}. A chain is displaced when the cascade frame differs from the event's natural frame. The displaced vs aligned impact comparison uses Mann-Whitney $U$. The reinforcement rate (\% catalyst) is tested per target frame with bootstrap 95\% CI (5,000 replicates).

\subsubsection*{Reproducibility}
Generated by \texttt{gen\_table\_triple\_chains.py} in \texttt{scripts/},
which reads StabSel Phase~1 and Phase~2 results.

